\documentclass[12pt]{article}
\usepackage[utf8]{inputenc}
\usepackage[T1]{fontenc}
\usepackage{hyperref}
\title{Pauserna Language Resources for K-12 Curriculum Development}
\author{Compiled by Claude Sonnet 5.5 (claude-sonnet-5-5) for Andr\'as Kornai}
\date{}
\begin{document}
\maketitle

\section{Introduction}
Pauserna (also Warázu, Guarasu'we, Guarasugwé; ISO 639-3: psm; Glottocode: paus1244) is a Tupian language of Bolivia and Brazil. Search summaries report only a handful of speakers (about 4--5 per Wikipedia, 2017; one speaker in Bolivia and a few in Brazil per a 2016 report by Danielsen), so it is critically endangered or dormant. Only general reference pages were found.

\section{Primary Educational Resources}
No primers, textbooks or school materials were found.

\section{Linguistic and Reference Materials}
\begin{itemize}
\item Warázu language (Wikipedia): \url{https://en.wikipedia.org/wiki/War%C3%A1zu_language}
\item Endangered Languages Project entry: \url{https://endangeredlanguages.com/lang/2957}
\item Verbix endangered languages page: \url{https://docs.verbix.com/EndangeredLanguages/Pauserna}
\end{itemize}
The summaries mention Danielsen and Sell's 2015 dictionary of Guarasu flora and fauna and Ramirez, Vegini and Fran\c{c}a's 2017 description (in Portuguese); no direct URLs for them were seen, so none is given.

\section{Religious and Community Resources}
\begin{itemize}
\item Encyclopedia entry on the Pauserna: \url{https://www.encyclopedia.com/humanities/encyclopedias-almanacs-transcripts-and-maps/pauserna}
\item EveryCulture entry on the Pauserna: \url{https://EveryCulture.com/South-America/Pauserna.html}
\end{itemize}

\section{Digital and Audio Resources}
No audio, video or software resources were found.

\section{Curriculum Development Context}
With so few speakers, a school programme would be revitalization from documentation, which exists mainly in academic work. The Bolivian community lives in the Territorio Guarasugwe-Chiquitano del Bajo Parag\'ua beside Chiquitano speakers, so any programme should be built with that community. Locating and licensing the published descriptions would be the first step.

\end{document}